% LAB BOREDOM — module L1, The Three Literatures. GUIDED DRAFTING (desktop recipe), Section 2 run.
% Session 2026-08-06. Author rulings this session: pack v2 is the bank; r1 ¶1 carried as working base.
% Pack: lab-boredom-L1-pack-v2 (+ -v2-quotes.json)  |  Spec: outline v7 §L1  |  Rulings: register v8.
% Protocol: Stage-3 guided batches (1–3 ¶¶, approval gate each batch), FCDP v2 conventions; «Qnn» markers,
%   mechanical substitution via scripts/substitute-quote-ids.py; rendered copy = L1-DRAFT-v1-rendered.tex.
% MASTER FILE — carries markers. Edit THIS file; regenerate the rendered copy after every edit.
%
% Corrections carried (tandem session, author-ruled): Q01 BARRED (abstract-only); Q02 @ printed 484;
%   Q05 attributed to Fahlman; O-06 settled = cite-not-re-present; E18 re-placed to ¶12.
% ¶1 carried verbatim from drafts/LAB-BOREDOM-L1-DRAFT-v1.tex batch-1 r1 (author-revised 2026-08-06).
%
% BATCH 1 = ¶¶1–3 (drafted 2026-08-06, awaiting author review).
% Batch-1 verification ledger:
%   Q14 printed page PINNED THIS SESSION: PDF p. 17 running head "The Other Side of Existence 117"
%     → printed p. 117 (folio-read, archon-cli-v3 corpus PDF). Slaby fn 18: Pessoa, Book of Disquiet,
%     fragment 219, p. 192. Pack v2 locus "RE-PIN" superseded by this pin; record in pack errata at batch close.
%   Q07 cited at VOLUME pp. 207–08 per pack (MLA cites the chapter in the volume).
%   Q12 cited at printed p. 101 per pack v2 locus pass.
% Batch-1 selection decisions (plan approved 2026-08-06):
%   Q08 HELD for ¶10 (its content paraphrased here, rendered there). Q09 PARAPHRASED (no volume folio
%   available; render would need a ****** pin). Q10 UNUSED this batch (optional add). Q11 PARAPHRASED-SKIP
%   (opener carries the chapter's question). Q13 not used (attribution hazard; not needed).
% Open in this batch:
%   ****** in ¶2 = DA-03 source names (Angst-vs-boredom rank dispute) — not in pack; pack-gap item.
%   E01 self-citing author unnamed (Stage-P resolution in flight) — ¶1 handles by not naming.

Boredom has been academically pursued across two rather different fields of study. One asks
what boredom discloses about the person who suffers it, treats the experience as a disclosure
of how things stand for someone rather than as a state to be relieved, and reaches for
Heidegger and the phenomenological tradition that follows him; the other asks what boredom
looks like in a signal, and reaches for electrodes, eye trackers, and self-report scales
validated against one another. Across the works surveyed here, the connection between them
amounts to a single author citing their own work across two programs, together with one
critical footnote. That is a citation trail rather than an exchange. The first counts a
reading of experience as evidence; the second counts a number extracted from a sensor.
Neither treats what the other counts as bearing on its own question, and the absence is not an
oversight either could have corrected by wider reading. The two studies we published stand in
neither literature: they belong to a third lineage, clinical immersion in engineering
education, where a headset substitutes for clinical access that no program can secure in
sufficient quantity, and that lineage arrives with its own question, its own warrant, its own
reference list, and no commitment about what boredom is. From the second literature we took
one thing, deliberately and by name: the definition of boredom under which we collected the
data. That definition reached us through Fahlman from Eastwood, it is the definition our
methods section declares, it governs what our instruments were built to detect, and it is also
the definition that could not name what we found. What can now be read out of this data
follows from that borrowing.

The phenomenological literature's liveliest quarrel concerns whether the experience Heidegger
analyzes is boredom at all. Elpidorou and Freeman set profound boredom against the construct
the empirical programs measure---a brief, situation-bound, many-sided response to
unsatisfying circumstances---and conclude of profound boredom that ``«Q07»'' (207--08). The
verdict protects both parties to the comparison: it spares the psychological construct a
metaphysics it never claimed, and it cautions every instrumented study, this one included,
against announcing that a sensor has caught what Heidegger meant. Nor has the reception
settled the analysis's rank; whether the fundamental attunement is profound boredom or
\textit{Angst} remains argued (******). Slaby reads affectivity as ``«Q12»'' (101): mood is
not noise laid over thinking but the way a person's standing in the world shows itself at
all. On that account a person can remain occupied---attentive, compliant, even performing
well---while what stands beneath the occupation is emptiness, the condition Pessoa's words
name most starkly---``«Q14»'' (qtd. in Slaby 117). An occupied surface over a hollow depth
is, under that reading, a structure description can reach, not a contradiction in terms.

Film is where the phenomenological literature comes nearest to this study's paradigm.
Quaranta reads film viewership through Heidegger's account of boredom, taking the watching of
a film as a mode of being-in-the-world in its own right; within the works surveyed here, his
is the one sustained reading of film through this frame. The claim matters for a plain
reason. The stimulus in this experiment is film: participants watched footage they did not
choose, under mandate, in a headset that held it before them. A literature prepared to treat
film viewership as a site where boredom's forms become visible is not being stretched to
reach this paradigm; the paradigm sits inside its subject matter.
